\documentclass[a4paper,12pt]{article}
\usepackage[utf8]{inputenc}
 \usepackage[english]{babel}
  \usepackage[backend=biber,natbib=true,maxnames=200]{biblatex}

 \addbibresource{AUTOSintro-refs.bib}

 \usepackage{graphicx}
 \usepackage{amssymb}
 \usepackage{amsmath}
 \usepackage{lscape}
 \usepackage{longtable}
 \usepackage{times}        % Benutzung des Times Zeichensatzes im Text
 \usepackage{caption2}     % Erweiterte Options für Bild- und Tabellenunterschriften
 \usepackage[pdftex]{hyperref}

 \setlength{\parskip}{12pt}
 \setlength{\parindent}{0pt}
 \setlength{\textwidth}{16.5cm}
 \setlength{\oddsidemargin}{0pt}
 \setlength{\evensidemargin}{0pt}
 \setlength{\unitlength}{1cm}

 \renewcommand\captionfont{\small}
 \renewcommand\captionlabelfont{\bfseries}
 \setcaptionwidth{\textwidth} % Breite der Bild- und Tabellenunterschriften
 \setlength{\LTcapwidth}{\textwidth} % Breite der longtable captions

 \newcommand{\tabletopspace}{\rule[0mm]{0mm}{5mm}} % legt den Abstand zwischen oberer
 % horizontaler Linie und erstem Eintrag fest

 \renewcommand{\arraystretch}{1.2}

 \renewcommand{\textfraction}{0.0}
 \renewcommand{\topfraction}{1.0}
 \renewcommand{\bottomfraction}{1.0}
 \newcommand{\threeJ}[6]{\ensuremath\left(\!\begin{array}{ccc} #1 & #2 & #3 \\ #4 & #5 & #6 \end{array}\!\right)}
 \newcommand{\sixJ}[6]{\ensuremath\left\{\!\begin{array}{ccc} #1 & #2 & #3 \\ #4 & #5 & #6 \end{array}\!\right\}}
 \newcommand{\CG}[6]{\ensuremath\langle #1 \, #2 ;\, #3 \, #4 \vert\, #5 \, #6 \rangle}
 \newcommand{\autos}{\textsc{Autostructure}}

\pagestyle{plain}
 %\includeonly{Abschnitt3}

\begin{document}

\textbf{\huge A short introduction to \autos}

\autos\ is a suite of atomic-structure codes that is particulary geared towards the calculation of recombination cross sections and rate coefficients. It has been coded in the programming language \textsc{Fortran} by Nigel \citet{Badnell2011b} from the University of Strathclyde (Glasgow, UK) and is available from a server there (\url{http://www.apap-network.org/autos/}). The code consists of the atomic-structure code \texttt{asdeck.f} and various post-processors for different type of data. For DR and RR calculations there are the post-processors \texttt{adasdr.f} and \texttt{adasrr.f}. Documentation is available as \url{http://www.apap-network.org/autos/ver/WRITEUP}.

\section{Installation}

A very basic installation procedure is described in the document \url{http://www.apap-network.org/autos/ver/INSTALL}. In order to be able to more conveniently handle the various input and output files that are read and generated by the various codes, a \texttt{makefile} and several shell scripts have been developed. Just type \texttt{"make"} to see all installation options. The essential individual installation steps are:
\begin{itemize}
\item compilation of \texttt{asdeck.f}: \texttt{"make AUTOS"},
\item compilation of \texttt{adasdr.f}: \texttt{"make AUTOSDR"},
\item to create soft-links in the \texttt{$\sim$/bin} directory to the shell scripts: \texttt{"make install"},
\item to make all shell scripts executable: \texttt{"chmod 755 *.sh"},
\item to extract documentation from the source files: \texttt{"make doc"},
\item to include the \texttt{$\sim$/bin} directory in your standard path: add the line\newline \texttt{"export PATH=\$PATH:\$HOME/bin"} to the file \texttt{.bashrc} in your home-directory.
\end{itemize}

\section{DR calculations with \autos}

\autos\ has to calculate all atomic level energies and transition rates which are required for the computation of DR cross sections. This is done partly by \texttt{asdeck.f} and partly (for high-$n$ levels) by \texttt{adasdr.f}. The  input files have to be named \texttt{<name>.in} and \texttt{<name>.drin}, respectively, where \texttt{<name>} can be chosen freely. All files pertaining to a given calculation are either stored in an automatically created subdirectory \texttt{<name>} or are named \texttt{<name>.<ext>} with various extensions \texttt{<ext>} as detailed below. In the following, DR of Be-like Ar$^{14+}$ is treated as an example.


\subsection{Level energies}\label{autos:DRlev}

Level energies of the Ar$^{14+}$ $1s^2\,2s^2$, $1s^2\,2s\,2p$, and $1s^2\,2p^2$ configurations.

\underline{Input file} \texttt{Ar14.in}:
\begin{verbatim}
A.S. Ar14+ energy levels
 &SALGEB CUP='IC' MXCONF=3 MXVORB=2 TITLE='Ar14'
         KCOR1=1 KCOR2=1 &END
 2 0   2 1
 2 0
 1 1
 0 2
 &SMINIM  NZION=-18 PRINT='FORM' &END
\end{verbatim}

Comments:
\begin{itemize}
\item \textsc{Fortran} input can be made via so-called namelists which group variables together. In the above example there are the two namelists \texttt{\&SALGEB} and \texttt{\&SMINIM}. The namelist input is terminated with \texttt{\&END}.
\item The ion charge is given as \texttt{NZION}. The negative sign refers to the computational method for the evaluation of radial wavefunctions (see \texttt{WRITEUP} for details).
\item The atomic configurations  are specified by the range of fixed orbitals. Presently, only the $1s$ orbital is fixed (\texttt{KCOR1=1, KCOR2=1}). In addition there are \texttt{MXVORB=2} orbitals, i.e., the $2s$ and $2p$ orbitals, which are specified as \texttt{"2 0~ 2 1"} on a separate line immediately below the namelist \texttt{\&SALGEB} input.
\item The subsequent  \texttt{MXCONF=3} lines specify the configurations \texttt{"2 0"} [$(1s^2)\,2s^2\,2p^0$],  \texttt{"1 1"} [$(1s^2)\,2s^1\,2p^1$], and  \texttt{"0 2"} [$(1s^2)\,2s^0\,2p^2$].
\item \texttt{"CUP='IC'"} selects the angular momentum coupling which is intermediate coupling in this case.
    \item \texttt{"PRINT='FORM'"} requests formatted human-readable output.
\end{itemize}

\underline{Executing the calculation}: Type \texttt{"autos Ar14"} on the command line.

\underline{Output file} \texttt{Ar14/LEVELS}:
\begin{verbatim}
2J P   S L   CF   NI       ENERGY(RYD)
 0 0   1 0    1    1        0.00000000
 0 1  -3 1    2    4        2.11717592
 2 1  -3 1    2    2        2.18220913
 4 1  -3 1    2    1        2.32660879
 2 1  -1 1    2    3        4.22709811
 0 0   3 1    3    4        5.56803877
 2 0   3 1    3    3        5.65434218
 4 0   3 1    3    2        5.77786583
 4 0   1 2    3    1        6.39944195
 0 0   1 0    3    5        7.78202330
 0 0   0 0    0    0     -758.69761466
\end{verbatim}
The level energies are in units of Rydbergs (1 Ryd = 13.606~eV).

\subsection{First step of DR calculation}\label{autos:DRstep1}

DR and TR of Ar$^{14+}$($1s^2\,2s^2$) via $2s\to2p$ and $2s^2\to2p^2$ excitations involving $1s^2\,2s\,2p\,n\ell$ and $1s^2\,2p^2\,n\ell$ intermediate configurations.

\underline{Input file} \texttt{Ar14DR.in}:
\begin{verbatim}
A.S. DR + TR of Ar14+
&SALGEB RUN='DR' RAD='YES' CUP='IC' MXCONF=3 MXCCF=3 MXVORB=2
       TITLE='2-2,n' KCOR1=1 KCOR2=1 LCON=5 &END
 2 0  2 1
 2 0
 1 1
 0 2

 2 1
 1 2
 0 3

 &DRR NMIN=3 NMAX=30 LMIN=0 LMAX=18 &END
 &SMINIM  NZION=-18 PRINT='UNFORM' &END
 &SRADCON MENG=-9  &END
  0.00  8.00
\end{verbatim}

Comments:
\begin{itemize}
\item Additional items are specified in the namelist \texttt{\&SALGEB}. There are \texttt{MXCFF=3} explicitly specified configurations of the recombined ion, i.e., \texttt{"2 1"} [$(1s^2)\,2s^2\,2p^1$],\newline  \texttt{"1 2"} [$(1s^2)\,2s^1\,2p^2$], and  \texttt{"0 3"} [$(1s^2)\,2s^0\,2p^3$]. These are the final $(1s^2)\,2s^2\,2p^1$ configuration and the intermediate $1s^2\,2p^2\,n\ell$ configurations with $n=2$.
\item The other $1s^2\,2p^2\,n\ell$ intermediate configurations with \texttt{NMIN} $\leq n \leq$ \texttt{NMAX} and \texttt{LMIN} $\leq \ell \leq$ \texttt{LMAX} are specified in the namelist \texttt{\&DRR}. Internally \autos\ extends the calculations to higher $n$ and $\ell$ values up to $n=999$ using hydrogenic rates and an interpolation scheme, where actually only a limited set of $n$-values is explicitly calculated. If you really want to stop the calculation at $n$ = \texttt{NMAX}, you need to add \texttt{NMESH} $=$ 0 to the namelist \texttt{\&DRR}.
\item The continuum energies of the initially free electron are specified in the namelist \texttt{\&SRADCON} and the line below. There are \texttt{$\vert$MENG$\vert$=9} energies in the energy range 0--8~Ryd. The negative sign of \texttt{MENG} is essential as explained in the \texttt{WRITEUP}. The highest energy should be larger than the largest excitation energy in the file \texttt{Ar14/LEVELS} calculated in section \ref{autos:DRlev}.
\end{itemize}

\underline{Executing the calculation}: Type \texttt{"autos Ar14DR"} on the command line.

\underline{Output files}: The file \texttt{Ar14DR/o1u} contains intermediate output which is to be processed in step 2 below. If something goes wrong, the file \texttt{Ar14DR/olg} may contain error messages.

\subsection{Second step of DR calculation}\label{autos:DRstep2}

\underline{Input file} \texttt{Ar14DR.drin}:
\begin{verbatim}
/IC/
&ONE COREX='2-2' NTAR1=2 &END
&TWO NR1=3 NBIN=4001 EMIN=0.0 EMAX=8.0
 EWIDTH=0.0 TPER=1.47e-3 TPAR=1.47e-5
 NECOR=10 &END
 0 0   1 0    1    1        0.00000000
 0 1  -3 1    2    4        2.11717592
 2 1  -3 1    2    2        2.18220913
 4 1  -3 1    2    1        2.32660879
 2 1  -1 1    2    3        4.22709811
 0 0   3 1    3    4        5.56803877
 2 0   3 1    3    3        5.65434218
 4 0   3 1    3    2        5.77786583
 4 0   1 2    3    1        6.39944195
 0 0   1 0    3    5        7.78202330
 0 0   0 0    0    0     -758.69761466
 0.0 2.0839 2.14931612 2.30258551 4.1206 5.51241 5.60557
   5.72556 6.39944195 7.6603
\end{verbatim}

Comments:
\begin{itemize}
\item Cross sections will be calculated for \texttt{"NBIN=4001"} electron energies ranging from \texttt{EMIN=0.0} to \texttt{EMAX=8.0}~Ryd.
\item \texttt{EWIDTH=0.0} flags the calculation of merged-beams rate coefficients (cf.~Eq.~\ref{eq:alphadelta}) with electron beam temperatures $k_BT_\perp=$ \texttt{TPER=1.47e-3}~Ryd = 20~meV and\newline $k_BT_\| = $ \texttt{TPAR=1.47e-5}~Ryd = 0.2~meV.
\item Results are produced for \texttt{NTAR1=2} initial levels, i.e., DR of Ar$^{14+}$($1s^2\,2s^2\;^1S_0$) and DR of Ar$^{14+}$($1s^2\,2s\,2p\;^3P_0$).
\item \texttt{NR1=3} specifies the lowest $n$ outside the core for which DR is calculated.
\item \texttt{COREX='2-2'} denotes that $\Delta n=0$ core excitation within the $L$-shell is considered.
\item The input below the namelist input is just a copy of the second to last line of the file \texttt{Ar14/LEVELS} calculated in section \ref{autos:DRlev}.
\item The very last line contains \texttt{NECOR=10} improved core-excitation energies (in Ryd, all on one line), e.g., from the NIST Atomic Spectra Database (\url{https://dx.doi.org/10.18434/T4W30F}).
\end{itemize}

\underline{Executing the calculation}: Type \texttt{"autosdr Ar14DR"} on the command line.

\underline{Output files}: The file \texttt{Ar14DR.adr} contains the cooler-convolved merged-beams rate coefficients (in $10^{-11}$~cm$^3$/s) as a function of electron energy (in Ryd). The file \texttt{Ar14DR.bdr} contains binned cross sections (in $10^{18}$~cm$^2$, see also below) which can be processed further, e.g., by the \textsc{hydrocal} code.

\newpage
\subsection{Binned cross sections}

The input file in section \ref{autos:DRstep2} defines an energy grid ranging from $E_\mathrm{min} = \texttt{EMIN}$ to $E_\mathrm{max} = \texttt{EMAX}$ containing $n_b = \texttt{NBIN}-1$ energy bins of width $\Delta E = (E_\mathrm{max}-E_\mathrm{min})/n_b$. Bin number $k$ ranges from $E^{(k)}_\mathrm{lo} = E_\mathrm{min}+k\cdot\Delta E$ to $E^{(k)}_\mathrm{hi} = E^{(k)}_\mathrm{lo}+\Delta E$. For a Lorentzian resonance with resonance energy $E_\mathrm{res}$, width $\Gamma$, and resonance strength $\overline{\sigma}$, i.e., for
\begin{equation}\label{eq:sigmaLor}
\sigma(E) = \frac{\overline{\sigma}} {\pi} \frac{\Gamma/2}{(E-E_\mathrm{res})^2+(\Gamma/2)^2}.
\end{equation}
the binned cross section for bin number $k$ is
\begin{equation}\label{eq:sigmabin1}
\sigma^{(k)} = \frac{\overline{\sigma}} {\pi\Delta E} \int_{E^{(k)}_\mathrm{lo}}^{E^{(k)}_\mathrm{hi}}\frac{\Gamma/2}{(E-E_\mathrm{res})^2+(\Gamma/2)^2}dE = \frac{\overline{\sigma}} {\pi\Delta E} \int_{\epsilon^{(k)}_\mathrm{lo}}^{\epsilon^{(k)}_\mathrm{hi}} \frac{1}{\epsilon^2+1}d\epsilon
\end{equation}
where
\[
\epsilon^{(k)}_\mathrm{lo} = \frac{E^{(k)}_\mathrm{lo}-E_\mathrm{res}}{\Gamma/2}\text{~~~and~~~}\epsilon^{(k)}_\mathrm{hi} = \frac{E^{(k)}_\mathrm{hi}-E_\mathrm{res}}{\Gamma/2}.
\]
The integral can be easily solved analytically yielding
\begin{equation}\label{eq:sigmabin2}
\sigma^{(k)} = \frac{\overline{\sigma}}{\pi\Delta E} \left[\arctan\left(\frac{E^{k}_\mathrm{hi}-E_\mathrm{res}}{\Gamma/2}\right) - \arctan\left(\frac{E^{k}_\mathrm{lo}-E_\mathrm{res}}{\Gamma/2}\right)\right].
\end{equation}
Internally \autos\ uses this type of binning \cite{Badnell2006b}.  The convolution with the merged-beams energy distribution  is carried out by representing the DR cross section as a sum of delta-like cross sections, i.e., by  approximating $\alpha_{MB}(\hat{E})$ as 
\begin{eqnarray}
\alpha_{MB}(\hat{E}) = \sum_k & & \!\!\!\!\!\!\!\frac{(\sigma^{(k)}\Delta E)\,c\sqrt{E^{(k)}}}{k_\mathrm{B}T_\perp\xi\sqrt{2m_ec^2}}\,\exp\left(-\frac{E^{(k)}-\hat{E}/\xi^2}{k_\mathrm{B}T_\perp}\right)\times\nonumber\\&&
  \left[\mathrm{erf}\left(\frac{\sqrt{E^{(k)}}+\sqrt{\hat{E}}/\xi^2}{\sqrt{k_\mathrm{B}T_\|}/\xi}\right)
             +\mathrm{erf}\left(\frac{\sqrt{E^{(k)}}-\sqrt{\hat{E}}/\xi^2}{\sqrt{k_\mathrm{B}T_\|}/\xi}\right) \right]
\end{eqnarray}
with $(\sigma^{(k)}\Delta E)$ from Eq.~\ref{eq:sigmabin2} and $E^{(k)} = (E^{k}_\mathrm{hi}+E^{k}_\mathrm{lo})/2 = E_\mathrm{min}+(k+0.5)\cdot\Delta E$

This approximation is sufficiently accurate as long as the bin width $\Delta E$ is much smaller than the experimental electron-energy spread \cite{Badnell2006b}. It should however be kept in mind that the representation of DR resonances as Lorentzians is not always appropriate, in particular, not at low energies. Moreover, the resonance energies are calculated as differences of large numbers which bears large uncertainties again for low energy resonances. All these limitations lead to noticeable differences between \autos\ results and experimental data, in particular, at low energies (see, e.g., \citep{Hahn2014a,Huang2018,Lestinsky2012,Lestinsky2009,Schmidt2008a,Wang2018b}).

\section{RR calculations with \autos}

\textbf{The recipe below produces some output, but its adequateness and its accuracy are currently still unclear. Please use with care!}


\autos\ has to calculate all atomic level energies and transition rates which are required for the computation of RR cross section. This is done partly by \texttt{asdeck.f} and partly (for high-$n$ levels) by \texttt{adasrr.f}. The  input files have to be named \texttt{<name>.in} and \texttt{<name>.rrin}, respectively, where \texttt{<name>} can be chosen freely. All files pertaining to a given calculation are either stored in an automatically created subdirectory \texttt{<name>} or are named \texttt{<name>.<ext>} with various extensions \texttt{<ext>} as detailed below. In the following, RR of O-like Ne$^{2+}$ is treated as an example.

\subsection{Levels energies}\label{autos:RRlev}

Level energies of the Ne$^{2+}$ $1s^2\,2s^2\,2p^4$, $1s^2\,2s^1\,2p^5$, and $1s^2\,2p^6$ configurations.

\underline{Input file:} \texttt{Ne2IC.in}
\begin{verbatim}
A.S. Ne2+
 &SALGEB RUN='  ' CUP='IC' MXCONF=3 MXVORB=2 TITLE='Ne2+'
         KCOR1=1 KCOR2=1 &END
 2 0  2 1
 2 4
 1 5
 0 6
 &SMINIM  NZION=-10 PRINT='FORM' &END
\end{verbatim}
See Section~\ref{autos:DRlev} for details.

\underline{Executing the calculation}: Type \texttt{"autos Ne2IC"} on the command line.

\underline{Output file} \texttt{Ne2IC/LEVELS}:
\begin{verbatim}
2J P   S L   CF   NI       ENERGY(RYD)
 4 0   3 1    1    2        0.00000000
 2 0   3 1    1    3        0.00609512
 0 0   3 1    1    4        0.00896683
 4 0   1 2    1    1        0.26198031
 0 0   1 0    1    5        0.50851848
 4 1  -3 1    2    1        2.01625991
 2 1  -3 1    2    2        2.02226537
 0 1  -3 1    2    4        2.02529861
 2 1  -1 1    2    3        2.91469946
 0 0   1 0    3    1        4.89144158
 0 0   0 0    0    0     -253.01521551
\end{verbatim}

The level energies are in units of Rydbergs (1 Ryd = 13.606~eV).

\subsection{First step of RR calculation}

RR of Ne$^{2+}$($1s^2\,2s^2\,2p^4$), capture into $nlj$-levels with $n\leq 15$.

\underline{Input file} \texttt{Ne2RR.in}:
\begin{verbatim}
A.S. Ne2+RR
 &SALGEB RUN='RR' RAD='E1' CUP='IC' MXVORB=2 MXCONF=3 TITLE='Ne2RR'
         KCOR1=1 KCOR2=1 LCON=3 &END
 2 0  2 1
 2 4
 1 5
 0 6
 &DRR NMIN=3 NMAX=15 LMIN=0 LMAX=14 &END
 &SMINIM  NZION=-10 PRINT='FORM' MAXE=10 MSTEP=12 &END
 &SRADCON EMAX=10 &END
\end{verbatim}

Basically a simplified input from Section~\ref{autos:DRstep1}. 
\begin{itemize}
\item Calculations are carried out for capture into $nl$ levels for $\texttt{NMIN} \leq n \leq \texttt{NMAX}$ and $\texttt{LMIN} \leq l\leq \texttt{LMAX}$. Higher $nl$ values are treated hydrogenically by the postprocessor (see Section \ref{autos:RRstep2} below).
\item Note that \texttt{MAXE} and \texttt{EMAX} are used in the namelists \texttt{SMINIM} and \texttt{SRADCON}, respectively, for specifying the maximum electron energy in Rydberg. 
\item \texttt{MSTEP} specifies the number of energy steps.
\end{itemize}

\underline{Executing the calculation}: Type \texttt{"autos Ne2RR"} on the command line.

\underline{Output files}: The file \texttt{Ne2RR/o1} contains intermediate output which is to be processed in step 2 below. If something goes wrong, the file \texttt{Ne2RR/olg} may contain error messages.

\subsection{Second step of RR calculation}\label{autos:RRstep2}

Summation of cross section. For high $n$ a hydrogenic approximation is used.

\underline{Input file} \texttt{Ne2RR.rrin}:
\begin{verbatim}
/IC/
&ONE NTAR1=1 NTAR2=10 IPRINT=-3 &END
&TWO NR1=3 JTHETA=0 &END
 4 0   3 1    1    2        0.00000000
 2 0   3 1    1    3        0.00609512
 0 0   3 1    1    4        0.00896683
 4 0   1 2    1    1        0.26198031
 0 0   1 0    1    5        0.50851848
 4 1  -3 1    2    1        2.01625991
 2 1  -3 1    2    2        2.02226537
 0 1  -3 1    2    4        2.02529861
 2 1  -1 1    2    3        2.91469946
 0 0   1 0    3    1        4.89144158
 0 0   0 0    0    0     -253.01521551
\end{verbatim}
Similar to input of Section \ref{autos:DRstep2}. Further comments:
\begin{itemize}\vspace*{-12pt}
\item \texttt{IPRINT=-3} requests minimal output.
\item \texttt{NR1=3} specifies the lowest $n$ outside the core for which DR is calculated.
\item \texttt{JTHETA=0} requests cross-section output.
\item The input below the namelist input is just a copy of the second to last line of the file \texttt{Ne2IC/LEVELS} calculated in section \ref{autos:RRlev}. It contains \texttt{NTAR2=10} levels.
\end{itemize}

\underline{Executing the calculation}: Type \texttt{"autosrr Ne2RR"} on the command line.

\underline{Output files}: The file \texttt{Ne2RR.arr} contains the RR cross section times the electron energy (in $10^{-18}$~cm$^2$ Ryd) as a function of electron energy (in Ryd).
\newpage
\printbibliography


\end{document}
